\documentclass[11pt]{article}

\usepackage[utf8]{inputenc}
\usepackage[T1]{fontenc}
\usepackage[margin=1in]{geometry}
\usepackage{amsmath,amssymb}
\usepackage{booktabs}
\usepackage{graphicx}
\usepackage{xcolor}
\usepackage[binary-units=true]{siunitx}
\usepackage[numbers,sort&compress]{natbib}
\usepackage[colorlinks=true,allcolors=blue]{hyperref}
\usepackage{authblk}

% ---------------------------------------------------------------------------
% Draft markers. Set \draftmodefalse to hide them before circulating.
% ---------------------------------------------------------------------------
\newif\ifdraftmode
\draftmodetrue

\ifdraftmode
  \newcommand{\todo}[1]{\textcolor{red}{\textbf{[TODO: #1]}}}
  \newcommand{\pending}{\textcolor{red}{\texttt{--.--}}}
  \newcommand{\note}[1]{\textcolor{blue!70!black}{\small\textsf{[#1]}}}
\else
  \newcommand{\todo}[1]{}
  \newcommand{\pending}{--}
  \newcommand{\note}[1]{}
\fi

\sisetup{detect-weight=true,detect-family=true}

\title{\bfseries Do learned molecular representations beat simple descriptors
for docking-score prediction?\\[0.3em]
\large A size-controlled benchmark at pool scale}

\author[1]{\todo{author list}}
\affil[1]{\todo{affiliation}}

\date{Draft --- \today}

\begin{document}
\maketitle

\begin{abstract}
Surrogate models that predict docking scores from molecular representations
underpin active-learning approaches to ultra-large virtual screening, where
exhaustive docking of billion-compound libraries is infeasible. A growing
number of pretrained molecular foundation models are used as the representation
in such surrogates, yet their reported performance is rarely compared against a
trivial descriptor baseline, and almost never controlled for molecular size ---
despite the fact that physics-based docking scores correlate strongly with the
number of heavy atoms. We benchmark \todo{N} pretrained molecular embeddings
against a \num{12}-feature physicochemical descriptor baseline for docking-score
regression on pools of up to \num{116} million molecules, under identical
splits, prediction head and evaluation. We report both conventional regression
and screening metrics and, critically, size-controlled variants computed by
partialling out heavy-atom count and by scoring within equal-count size strata.
We find that the descriptor baseline reaches a Pearson correlation of
\num{0.891} on the EnamineHTS/DOCK3.7 pool, saturating after roughly
\num{200000} training examples, and that this performance is \emph{not} a size
artifact: it retains a within-stratum correlation of \num{0.800} and an
enrichment factor of \num{30.1} at the top \SI{1}{\percent}.
\todo{headline claim about learned embeddings --- pending sweep}
We release the full pipeline, including the dataset alignment checks required
to reproduce the benchmark.
\end{abstract}

\section{Introduction}
\label{sec:intro}

Make-on-demand chemical libraries now exceed \num{e10} compounds, far beyond
what structure-based docking can exhaustively evaluate. Ultra-large docking
campaigns have shown that screening deeper into such libraries recovers
better and more diverse chemotypes~\citep{lyu2019ultralarge}, but the cost of
docking every molecule grows with library size. The dominant response is to
dock a subset and train a \emph{surrogate} model that predicts the docking
score of the remainder from a cheap molecular representation, then dock only
the molecules the surrogate ranks highly. This is the basis of pool-based
active learning for virtual screening~\citep{graff2021molpal} and of related
iterative schemes~\citep{gentile2022deepdocking}.

The surrogate's representation is therefore a consequential design choice, and
it has become common to use a pretrained molecular foundation model for it.
Such models are available in increasing number and variety --- trained on
SMILES~\citep{ross2022molformer,soares2025smited}, on
SELFIES, on molecular graphs~\citep{rong2020grover}, on 3D
conformers~\citep{zhou2023unimol}, and on 2D depictions --- and recent work has
systematically benchmarked how combining them affects downstream prediction
tasks in drug discovery~\citep{martinezgalindo2026multimodality}. That study
trained over \num{1400} models across \num{29} benchmarks and concluded that
additional modalities generally help provided an effective fusion strategy is
used, with late fusion outperforming intermediate fusion.

Two gaps motivate the present work.

\paragraph{The missing baseline.}
Benchmarks of learned molecular representations are typically compared against
each other, or against a fingerprint, rather than against cheap
physicochemical descriptors. For docking-score prediction this omission is
consequential, because docking scoring functions are not scale-free: a
DOCK3.7 score is a sum of per-atom interaction terms and correlates strongly
with molecular size. A representation that encodes little beyond size can
therefore post a high correlation against the true score while carrying no
information about binding. Without a descriptor control, it is not possible to
tell these cases apart.

\paragraph{The missing size control.}
The same mechanism contaminates the evaluation metrics themselves. Pearson
correlation and RMSE computed over a whole pool are dominated by the
size-driven component of the variance. A screener, however, operates within a
desired physicochemical envelope and needs a model that discriminates
\emph{among} molecules of comparable size. Metrics that do not control for size
do not answer that question.

\paragraph{Contributions.}
\begin{enumerate}
  \item We establish a \num{12}-feature physicochemical descriptor baseline for
        docking-score regression, evaluated under identical splits, prediction
        head and metrics as the learned representations it is compared
        against (Section~\ref{sec:methods}).
  \item We introduce a size-controlled evaluation protocol --- partial
        correlation, residual correlation after a polynomial size fit, and
        metrics within equal-count heavy-atom strata --- and show it separates
        genuine chemical signal from size encoding
        (Section~\ref{sec:results-size}).
  \item We benchmark \todo{N} pretrained molecular embeddings, spanning
        \numrange{512}{3400} dimensions, against that baseline on pools of up to
        \num{116} million molecules, reporting learning curves rather than
        single operating points (Section~\ref{sec:results-backbones}).
  \item We examine late fusion across representations and show that the
        prediction-averaging strategy reported as effective in the multimodal
        benchmarking literature degrades sharply when member quality is uneven
        (Section~\ref{sec:results-fusion}).
  \item We release a fully reproducible pipeline, including the dataset
        alignment verification that this benchmark turns out to require
        (Section~\ref{sec:repro}).
\end{enumerate}

\todo{Finalise the headline framing once the backbone sweep completes. If
learned embeddings fail to clear the descriptor baseline, the paper is a
negative result about representation quality and the title should say so. If
they clear it comfortably, the contribution is the size-controlled protocol
plus a cost--accuracy analysis, and the framing shifts accordingly.}

\section{Related work}
\label{sec:related}

\paragraph{Ultra-large virtual screening.}
Docking campaigns against make-on-demand libraries have recovered potent
ligands at hit rates that improve with library
size~\citep{lyu2019ultralarge}. The AmpC $\beta$-lactamase and D4 dopamine
receptor screens from that work are the source of two of the pools used here,
and their public release of docked scores is what makes surrogate benchmarking
on realistic distributions possible at all.

\paragraph{Surrogate models and active learning.}
MolPAL framed accelerated screening as pool-based active
learning, docking an initial batch, fitting a surrogate, and acquiring further
batches by predicted score or upper confidence
bound~\citep{graff2021molpal}. Deep Docking applies a related iterative
filtering scheme~\citep{gentile2022deepdocking}. Both report that a
surrogate trained on a small fraction of the pool recovers a large fraction of
the top-scoring molecules. Their representations are typically fingerprints or
descriptors; our question is what pretrained embeddings add relative to that
starting point.

\paragraph{Molecular representation learning.}
Pretrained molecular models differ in the modality they observe: SMILES
sequence models~\citep{ross2022molformer,soares2025smited}, graph-based
models~\citep{rong2020grover}, hypergraph-grammar
autoencoders~\citep{kajino2019mhg}, and 3D conformer
models~\citep{zhou2023unimol}. Extended-connectivity
fingerprints~\citep{rogers2010ecfp} remain a strong classical baseline. Each
carries a distinct inductive bias, and a recurring question is whether those
biases are complementary.

\paragraph{Multimodal fusion for molecular property prediction.}
\citet{martinezgalindo2026multimodality} conducted the most systematic study to
date, training over \num{1400} models across \num{7} molecular modalities,
\num{2} architectures, \num{2} fusion strategies and \num{29} benchmarks drawn
from the Therapeutics Data Commons~\citep{huang2021tdc}. They report that
additional modalities improve average performance, that late fusion
consistently outperforms intermediate fusion, and that knowledge-enhanced
pretraining yields further gains. Our work differs in three respects. First,
the target is a physics-based docking score rather than an experimental
property, so the label is cheap, abundant, and --- crucially --- has a known
dependence on molecular size that permits an explicit control. Second, the
pools are three to five orders of magnitude larger, which turns representation
dimensionality into a storage and compute decision rather than an academic one.
Third, we include a descriptor control, which that study and most of the
benchmarks it builds on do not.

\todo{Verify whether a size-controlled descriptor baseline for docking
surrogates has already been published. If so, reframe this contribution as
quantification at scale rather than as a novel control. Candidate literature:
deep-docking follow-ups, MolPAL ablations, QSAR-era work on the size dependence
of scoring functions.}

\section{Data}
\label{sec:data}

\subsection{Pools}

Each pool pairs a single receptor with a molecular library scored by one
docking program, so every row is one molecule and the label is its docking
score (more negative indicates a better predicted binder). Because the receptor
is fixed within a pool, no protein representation enters the model: the task is
a regression on molecular representation alone.

\begin{table}[t]
\centering
\caption{Docking pools. Scores are not comparable in absolute terms across
oracles; recall and enrichment are.}
\label{tab:pools}
\small
\begin{tabular}{llrrl}
\toprule
Pool & Target & Molecules & Backbones & Oracle \\
\midrule
\texttt{enhits} & EnamineHTS                & \num{2104318}    & 5 & DOCK3.7 \\
\texttt{ampc}   & AmpC $\beta$-lactamase    & \num{98489350}   & 6 & DOCK3.7 \\
\texttt{d4}     & D4 dopamine receptor      & \num{116225927}  & 6 & DOCK3.7 \\
\texttt{sura}   & SurA chaperone            & \num{1531875935} & \todo{--} & Glide HTVS \\
\bottomrule
\end{tabular}
\end{table}

The \texttt{sura} pool is listed for completeness but excluded from all
experiments reported here: its embeddings exist only as partial, unstitched
per-chunk arrays at the time of writing. Its inclusion matters for a future
version of this work precisely because it uses a different docking program;
until then, every result below is a statement about DOCK3.7 specifically rather
than about docking in general.

\subsection{Representations}

Each pool is accompanied by precomputed embeddings from several pretrained
molecular models, stored as dense \texttt{float32} arrays row-aligned with the
molecule list (Table~\ref{tab:backbones}). Storage is a material constraint: a
single \num{3400}-dimensional embedding of the D4 pool occupies
\SI{1.58}{\tera\byte}.

\begin{table}[t]
\centering
\caption{Embedding backbones, dimensionality, and storage for the D4 pool
(\num{116225927} molecules, \texttt{float32}).}
\label{tab:backbones}
\small
\begin{tabular}{llrr}
\toprule
Backbone & Modality & Dim. & D4 storage (\si{\tera\byte}) \\
\midrule
GROVER~\citep{rong2020grover}        & molecular graph     & 3400 & 1.58 \\
Uni-Mol2~\citep{zhou2023unimol}      & 3D conformer        & 1536 & 0.71 \\
MHG-GED~\citep{kajino2019mhg}        & hypergraph grammar  & 1024 & 0.48 \\
MolFormer~\citep{ross2022molformer}  & SMILES              &  768 & 0.36 \\
SMI-TED~\citep{soares2025smited}     & SMILES              &  768 & 0.36 \\
Uni-Mol v1~\citep{zhou2023unimol}    & 3D conformer        &  512 & 0.24 \\
\bottomrule
\end{tabular}
\end{table}

Two inconsistencies in the stored artifacts are worth recording, as they are
invisible to a consumer who assumes uniformity. On \texttt{enhits}, the
Uni-Mol2 embedding is \num{768}-dimensional rather than \num{1536}, having been
produced by a different extraction; results for that backbone are consequently
not comparable across pools. The same pool also lacks Uni-Mol v1 entirely and
stores GROVER under a different filename convention.

\subsection{Row alignment}
\label{sec:alignment}

Surrogate benchmarking of this kind depends on an assumption that is easy to
state and easy to violate: that row $i$ of the embedding array corresponds to
row $i$ of the score vector. A violation does not raise an error. It produces a
model that trains to a plausible-looking loss curve and reports meaningless
numbers.

For two of the three pools used here the assumption holds directly. For
\texttt{enhits} it does not: the embeddings are stored in a shuffled order
relative to the only available score file, which is itself sorted by score, so
a naive join yields labels that are both permuted and systematically ordered.
We therefore treat alignment as a verification step rather than an assumption.
Our pipeline reconstructs the permutation, applies it, and then checks the
molecular identifier of every row against a reference list, refusing to proceed
on any mismatch. All \num{2104318} rows of \texttt{enhits} pass this check; we
additionally confirmed agreement on \num{50000} randomly drawn rows through an
independent identifier-to-score mapping.

We report this in detail because it is a reproducibility hazard rather than an
incidental engineering note: any benchmark built on these public artifacts
without such a check may silently be reporting noise.

\subsection{Failed docking runs}

Docking runs that fail to produce a pose, or that produce severe steric
clashes, yield non-negative scores that are not on the same scale as successful
results. These are removed before splitting, so that training, validation and
test sets share one universe and the filter is recorded alongside the split
rather than applied ad hoc at training time. The affected fraction varies
sharply by pool: approximately \SI{5}{\percent} of the D4 pool (maximum
observed score $+328$), a negligible fraction of AmpC, and \num{16} rows of
\texttt{enhits}.

\section{Methods}
\label{sec:methods}

\subsection{Splits}

For each pool we draw a single random partition into training, validation and
test sets (\SI{90}{\percent}/\SI{5}{\percent}/\SI{5}{\percent}) after removing
failed docking runs. The same partition is reused by every backbone, so a
comparison between representations is never confounded by a difference in the
underlying rows. Training indices are stored in random order, which allows a
training budget of $N$ to be realised as the first $N$ indices; learning-curve
points are therefore nested subsets rather than independent samples, and the
corresponding embedding cache for budget $N$ is a prefix of the full cache.

Model selection uses the validation set only. The test set is read once per
run, after the best epoch has been selected, and the resulting predictions are
written to disk. This last detail is what makes the fusion and size-control
analyses in Sections~\ref{sec:results-fusion} and~\ref{sec:results-size}
possible without retraining.

\subsection{Prediction head}

To isolate the effect of the representation, the prediction head is held fixed
across all backbones: a multilayer perceptron with hidden widths
$(1024, 512)$, batch normalisation, ReLU activations and dropout $0.1$,
regressing a standardised score under mean squared error. Inputs are
standardised per feature using statistics estimated on a training subsample.
We optimise with AdamW (learning rate \num{e-3}, weight decay \num{e-5}) under
a cosine schedule, for at most \num{40} epochs with early stopping after
\num{6} epochs without validation improvement.

Holding the head fixed trades per-backbone optimality for comparability. We
address the resulting objection --- that a single architecture may suit some
representations better than others --- in Section~\ref{sec:results-robust}.

\subsection{Descriptor baseline}
\label{sec:methods-baseline}

The baseline replaces the learned embedding with \num{12} physicochemical
descriptors computed with RDKit~\citep{rdkit}: heavy atom count, molecular
weight, \mbox{cLogP}, topological polar surface area, ring count, rotatable
bond count, hydrogen bond donor and acceptor counts, fraction of sp\textsuperscript{3}
carbons, halogen count, aromatic ring count, and heteroatom count. These pass
through the identical split, prediction head, optimiser and metrics as every
learned representation, so the comparison is of representation alone.

\subsection{Metrics}

Lower scores indicate better predicted binders, so ``top-$k$'' refers to the
$k$ most negative values. We report regression metrics (RMSE, MAE, $R^2$,
Pearson $r$, Spearman $\rho$) and screening metrics: recall of the true top
fraction $f$ within the predicted top fraction $f$, and the corresponding
enrichment factor $\mathrm{EF}@f = \mathrm{recall}@f / f$, at
$f \in \{0.001, 0.01, 0.05\}$.

The two families answer different questions and need not agree. Pearson $r$
and RMSE are computed over the whole pool and are therefore dominated by the
large majority of non-binders; enrichment at \SI{0.1}{\percent} describes the
only region of the ranking a screening campaign would act upon.

\subsection{Size-controlled evaluation}
\label{sec:methods-size}

Because docking scores scale with molecular size, we compute three additional
quantities for every set of predictions, all from the stored test predictions
and therefore applicable retrospectively.

\paragraph{Partial correlation.}
The correlation between true and predicted score with heavy-atom count $s$
partialled out of both,
\begin{equation}
r_{yp \cdot s} = \frac{r_{yp} - r_{ys} r_{ps}}
                      {\sqrt{(1 - r_{ys}^2)(1 - r_{ps}^2)}},
\end{equation}
which removes a linear dependence on size.

\paragraph{Residual correlation.}
The correlation between the residuals of $y$ and of $\hat{y}$ after each is
regressed on a degree-3 polynomial in $s$, which permits size to act
non-linearly.

\paragraph{Within-stratum metrics.}
The test set is partitioned into ten equal-count bins of heavy-atom count, and
Pearson $r$, Spearman $\rho$ and $\mathrm{EF}@0.01$ are computed inside each
bin and averaged with weights proportional to bin size. This assumes no
functional form for the size dependence, and corresponds to the practical
question of whether a model can rank molecules of comparable size.

As a reference point we also evaluate the degenerate predictor
$\hat{y} = -s$, which by construction attains zero partial and residual
correlation and near-chance within-stratum performance. Its raw correlation
quantifies how much of a pool's score variance size alone explains.

\subsection{Fusion}

We combine representations by late fusion: each backbone is trained
independently and their test predictions are averaged. This permits any subset
to be evaluated without retraining, which is the property that made late fusion
attractive in prior multimodal
benchmarks~\citep{martinezgalindo2026multimodality}. We additionally consider
rank averaging and standardised averaging, which discard calibration in
exchange for robustness to differences in scale between members.
\todo{Add intermediate fusion (concatenated embeddings, jointly trained) if the
multimodality framing is retained.}

\section{Results}
\label{sec:results}

Unless stated otherwise, results are for the \texttt{enhits} pool
(EnamineHTS, DOCK3.7), with a test set of \num{105215} molecules held out from
\num{2104302} usable rows, under a single random seed.
\todo{Repeat all headline comparisons over five seeds and report dispersion;
single-seed numbers are provisional throughout.}

\subsection{The descriptor baseline and its learning curve}
\label{sec:results-baseline}

Table~\ref{tab:baseline-curve} reports the descriptor baseline as a function of
training set size. Performance saturates early: increasing the training set by
a factor of \num{38}, from \num{50000} to \num{1893872} molecules, improves
Pearson $r$ by \num{0.028}. This is the expected behaviour of a \num{12}-feature
model, which exhausts its capacity long before it exhausts the available data.

\begin{table}[t]
\centering
\caption{Descriptor baseline on \texttt{enhits}, as a function of training set
size. Test set \num{105215} molecules.}
\label{tab:baseline-curve}
\small
\begin{tabular}{rrrrrr}
\toprule
Train size & Pearson $r$ & Spearman $\rho$ & RMSE & $\mathrm{EF}@1\%$ & $\mathrm{EF}@0.1\%$ \\
\midrule
\num{50000}   & 0.8748 & 0.8613 & 0.4611 & 37.5 & 238 \\
\num{200000}  & 0.8802 & 0.8669 & 0.4517 & 38.4 & 248 \\
\num{1000000} & 0.8872 & 0.8744 & 0.4391 & 38.7 & 257 \\
\num{1893872} & 0.8907 & 0.8783 & 0.4326 & 40.7 & 286 \\
\bottomrule
\end{tabular}
\end{table}

The practical consequence is that the bar a learned representation must clear
is not only high but cheap: a Pearson $r$ of \num{0.88} is available from
twelve descriptors and \num{200000} docked molecules, with no embedding
storage and no pretraining.

\subsection{The baseline is not an artifact of molecular size}
\label{sec:results-size}

Heavy-atom count alone correlates $-0.659$ with the DOCK3.7 score on this pool,
which raises the possibility that the baseline's apparent strength is simply
size encoding. Table~\ref{tab:sizecontrol} shows that it is not.

\begin{table}[t]
\centering
\caption{Size-controlled evaluation on \texttt{enhits}. ``Raw'' is the
conventional metric; ``partial'' and ``residual'' remove a linear and a
degree-3 dependence on heavy-atom count respectively; ``within-stratum''
quantities are computed inside ten equal-count heavy-atom bins and weighted by
bin size. The first row is the degenerate size-only predictor.}
\label{tab:sizecontrol}
\small
\begin{tabular}{lrrrrr}
\toprule
 & \multicolumn{3}{c}{Pearson $r$} & \multicolumn{2}{c}{Within-stratum} \\
\cmidrule(lr){2-4}\cmidrule(lr){5-6}
Model & raw & partial & residual & $r$ & $\mathrm{EF}@1\%$ \\
\midrule
Size only ($-$heavy atoms) & 0.6592 & 0.000 & 0.000 & 0.1478 & 2.6 \\
\midrule
Descriptors, \num{50000}   & 0.8748 & 0.7647 & 0.7598 & 0.7692 & 25.8 \\
Descriptors, \num{200000}  & 0.8802 & 0.7756 & 0.7710 & 0.7797 & 28.1 \\
Descriptors, \num{1000000} & 0.8872 & 0.7896 & 0.7853 & 0.7933 & 29.6 \\
Descriptors, \num{1893872} & 0.8907 & 0.7966 & 0.7924 & 0.8000 & 30.1 \\
\bottomrule
\end{tabular}
\end{table}

Two observations follow. First, the control behaves as intended: the size-only
predictor attains a substantial raw correlation of \num{0.659} but collapses to
a within-stratum correlation of \num{0.148} and an enrichment factor of
\num{2.6}, close to the chance value of \num{1.0}. A feature that is
near-constant within its own stratum cannot discriminate inside it.

Second, the descriptor baseline loses comparatively little. At the largest
training size its correlation falls from \num{0.891} to \num{0.797} under
partial correlation and \num{0.800} within strata, and its top-\SI{1}{\percent}
enrichment from \num{40.7} to \num{30.1}. Approximately \SI{26}{\percent} of
the raw enrichment is attributable to size; the remainder reflects genuine
chemical information carried by the non-size descriptors.

This result is important for the interpretation of everything that follows. A
descriptor baseline that was merely a size proxy would be a weak and somewhat
unfair comparator. One that retains a within-stratum enrichment of \num{30}
is a substantive target, and failure to exceed it is correspondingly
informative.

\subsection{Learned representations}
\label{sec:results-backbones}

\todo{Primary results table pending completion of the backbone sweep. Populate
with: backbone, dimensionality, storage, raw and size-controlled metrics, at
each of four training budgets. Expected structure below.}

\begin{table}[t]
\centering
\caption{Learned representations versus the descriptor baseline on
\texttt{enhits} at the largest training budget. \todo{pending}}
\label{tab:backbones-results}
\small
\begin{tabular}{lrrrrrr}
\toprule
Representation & Dim. & Pearson $r$ & partial $r$ & within $r$ & $\mathrm{EF}@1\%$ & within $\mathrm{EF}@1\%$ \\
\midrule
Descriptors & 12 & 0.8907 & 0.7966 & 0.8000 & 40.7 & 30.1 \\
\midrule
MHG-GED     & 1024 & \pending & \pending & \pending & \pending & \pending \\
GROVER      & 3400 & \pending & \pending & \pending & \pending & \pending \\
Uni-Mol2    &  768 & \pending & \pending & \pending & \pending & \pending \\
MolFormer   &  768 & \pending & \pending & \pending & \pending & \pending \\
SMI-TED     &  768 & \pending & \pending & \pending & \pending & \pending \\
\bottomrule
\end{tabular}
\end{table}

\todo{Learning-curve figure: Pearson $r$ against training size, one line per
representation, with the descriptor baseline drawn as a horizontal reference.
The scientifically informative feature is curve \emph{shape}: a representation
that begins below the baseline but continues to improve where the baseline has
flattened has capacity that additional docking would unlock, which is a
different conclusion from one that flattens at the same point.}

\subsection{Robustness to the choice of prediction head}
\label{sec:results-robust}

A fixed prediction head invites the objection that any observed deficit of
learned representations reflects an architecture tuned, in effect, for
low-dimensional descriptors. We therefore repeat the comparison with two
additional model families on the same representations and splits: ridge
regression, and gradient-boosted decision trees.

\todo{Run and report. If the ordering between representations is preserved
across three model families of markedly different inductive bias, the result
concerns the representation rather than the head. This is the single most
important robustness check in the paper.}

\subsection{Late fusion}
\label{sec:results-fusion}

\todo{Rerun fusion on converged models; the numbers below are from an
undertrained preliminary sweep and are reported only to indicate the direction
of the effect.}

In a preliminary experiment, averaging the predictions of a strong
representation with two weaker ones reduced Pearson $r$ from \num{0.938} to
\num{0.839}. Prediction averaging weights all members equally, so the
combination is pulled toward its weakest constituent. Prior reports that late
fusion reliably improves on single
modalities~\citep{martinezgalindo2026multimodality} were obtained in a regime
where member representations performed comparably; that condition does not hold
here, and unweighted fusion is correspondingly harmful.

\todo{Quantify under the converged sweep, and compare unweighted averaging
against rank averaging and validation-weighted averaging. A gating rule ---
fuse only members that clear the descriptor baseline --- is the natural
remedy and should be evaluated.}

\subsection{Transfer across targets}

\todo{Repeat the headline comparison on \texttt{ampc} and \texttt{d4}. Both
use DOCK3.7, so agreement across them establishes robustness to target but not
to scoring function. Claims about docking in general require the Glide pool,
which is presently unavailable (Section~\ref{sec:data}).}

\section{Discussion}
\label{sec:discussion}

\todo{Write once the backbone sweep and the robustness check are complete. The
argument differs substantially between outcomes; sketches of both are retained
below so that the framing is a deliberate choice rather than a reaction to the
numbers.}

\paragraph{If learned representations do not clear the baseline.}
The claim would be that pretrained molecular embeddings, as commonly used,
carry little information about docking score beyond what twelve
physicochemical descriptors already provide --- and that reported successes
may partly reflect the absence of such a control together with metrics that do
not account for molecular size. The practical implication is immediate: a
surrogate for ultra-large screening need not store terabytes of embeddings,
and the storage freed is better spent docking more molecules. The claim would
require the robustness check of Section~\ref{sec:results-robust} to hold across
model families, and replication on at least one further target.

\paragraph{If they clear it comfortably.}
The contributions become the size-controlled protocol itself, the demonstration
that a substantial share of apparent performance is size-driven, and a
cost--accuracy analysis over embedding dimensionality. The relevant question
shifts to whether the highest-dimensional representations justify their
storage: GROVER occupies \num{6.6} times the bytes per molecule of Uni-Mol v1,
and a comparable accuracy at that ratio is a decision-relevant finding for
anyone provisioning a screening campaign.

\paragraph{In either case.}
Two observations appear robust to the outcome. First, the descriptor baseline
is strong, cheap and non-trivial once size is controlled for, and its absence
from the surrogate literature is a gap worth closing regardless of which side
of it learned representations fall. Second, unweighted late fusion is unsafe
when member quality is uneven; conclusions drawn from benchmarks of comparable
modalities do not transfer automatically to settings where some
representations are much weaker than others.

\section{Limitations}
\label{sec:limitations}

\paragraph{One scoring function.}
All results reported here come from DOCK3.7. The size dependence that motivates
our control is a property of the scoring function, and its magnitude will
differ for other programs. Until the Glide pool is available, claims should be
read as statements about DOCK3.7 rather than about docking generally.

\paragraph{Docking scores are not binding affinities.}
A surrogate that reproduces a docking score reproduces the scoring function's
errors along with its signal. Strong performance here implies nothing about
experimental potency, and the enrichment metrics measure recovery of
top-scoring molecules, not of actual binders.

\paragraph{Random splits.}
Our splits are random, which matches the deployment setting of a surrogate
applied within one pool, but gives no information about generalisation to
novel scaffolds. Scaffold-based or cluster-based splits would answer a
different and also important question.

\paragraph{A single prediction head in the primary comparison.}
The main results hold the head fixed for comparability.
Section~\ref{sec:results-robust} is the mitigation; absent that check, the
comparison is between (representation, head) pairs rather than between
representations.

\paragraph{Embeddings taken as given.}
We consume precomputed embeddings and do not fine-tune the underlying models.
Fine-tuning on docking data might substantially change the ordering, and our
results speak only to the frozen-representation setting --- which is, however,
the setting in which these models are typically deployed at pool scale, since
fine-tuning across a billion molecules is itself prohibitive.

\paragraph{Nested learning-curve subsets.}
Training budgets are nested prefixes of one shuffled index list rather than
independently drawn samples, so points along a learning curve are correlated.
This reduces variance in the comparison between budgets but does not support
treating them as independent observations.

\todo{Add statistical power statement once multi-seed runs are complete.}

\section{Reproducibility}
\label{sec:repro}

Code for the full pipeline --- score preparation with identifier verification,
split construction, training, fusion, reporting and size-controlled evaluation
--- is released at \todo{repository URL}. The docking scores and molecular
libraries are public~\citep{lyu2019ultralarge,graff2021molpal}.

We draw attention to one aspect in particular. As described in
Section~\ref{sec:alignment}, reproducing this benchmark requires reconstructing
a permutation between the stored embedding order and the stored score order for
one of the pools, and the failure mode when this is omitted is silent: training
proceeds normally and the resulting metrics are meaningless. Our preparation
step therefore verifies the molecular identifier of every row against a
reference list and refuses to produce a score array on any mismatch. We
recommend that benchmarks built on precomputed embedding artifacts adopt an
equivalent check, and we report ours explicitly so that discrepancies with
future work can be localised.

\todo{State final software versions, hardware, and total compute. Record the
exact commit hash used for the reported results.}


\bibliographystyle{unsrtnat}
\bibliography{refs}

\end{document}
